\section{Regularization transported by the output channel}\label{sec:ridge84}
The fixed Euclidean ridge in Theorem~\ref{thm:closedcovariance83}
is tied to the displayed outcome coordinates.  Its contraction under
arbitrary stochastic relabeling does not follow from the
unregularized data-processing theorem.  A ridge obtained from a
public probability vector has a different, covariant property.

For a probability vector $w=(w_j)$, put $W=(w_jI_d)_j$ and
$\mathcal B_w=\mathcal C_W$.  On the full real Hermitian tuple
space this operator is
\begin{equation}\label{eq:weightedridge84}
 (\mathcal B_wK)_j=w_j\left(K_j-\sum_\ell w_\ell K_\ell\right).
\end{equation}
For a zero-sum direction $H$ and $\tau>0$ define the extended energy
\begin{equation}\label{eq:transportenergy84}
 \mathcal E_{E,w,\tau}(H)=\sup_{K\in\mathcal H_d^k}
 \left\{2\langle H,K\rangle-
       \langle K,(\mathcal C_E+\tau\mathcal B_w)K\rangle\right\}.
\end{equation}
It equals the pseudoinverse quadratic form on the range and is
$+\infty$ off the range.  Constant tuples have zero energy and
zero pairing with $H$.  If every $w_j>0$, the supremum is finite
and may be evaluated by restricting to $\mathcal T$.

\begin{theorem}[Processing-covariant regularized energy]\label{thm:transportedridge84}
For every column-stochastic output map $T$, every $\tau>0$, and
every zero-sum Hermitian tuple $H$,
\begin{equation}\label{eq:ridgeprocessing84}
 \mathcal E_{TE,Tw,\tau}(TH)\le\mathcal E_{E,w,\tau}(H).
\end{equation}
Zero entries of $Tw$ are allowed, with the extended-energy convention
in \eqref{eq:transportenergy84}.  If a stochastic map $S$ satisfies
$ST=\operatorname{Id}$, equality holds in
\eqref{eq:ridgeprocessing84}.  In particular, reversible outcome
splitting leaves this energy unchanged when the public weights
are split by the same rule.
\end{theorem}
\begin{proof}
The matrix Jensen identity in
Proposition~\ref{prop:covarianceprocessing83} is an inequality on
the full tuple space, not only on its zero-sum subspace.  It gives
\[
 \langle L,\mathcal C_{TE}L\rangle
 \ge\langle T^*L,\mathcal C_E T^*L\rangle.
\]
Apply the same inequality to the scalar measurement $W$ to get
$\langle L,\mathcal B_{Tw}L\rangle
\ge\langle T^*L,\mathcal B_wT^*L\rangle$.
Also $\langle TH,L\rangle=\langle H,T^*L\rangle$.
Thus each expression in the supremum defining the left side of
\eqref{eq:ridgeprocessing84} is bounded by a permissible expression
on the right.  Taking suprema proves the claim, including infinite
values.  No assumption that $T^*L$ has zero sum is needed.
If $ST=\operatorname{Id}$, application of the same inequality to
$S$ gives the reverse bound.
\end{proof}

\begin{proposition}[Relation to the operational upper certificate]\label{prop:ridgeupper84}
If $w_j>0$ and $w_{\max}=\max_jw_j$, then for every pair of legal
measurements, with midpoint $M$ and difference $H$,
\begin{equation}\label{eq:ridgeupper84}
 D_N(E,F)\le\min\left\{2,
 2\sqrt{k+1}\sqrt{N\mathcal E_{M,w,1/(Nw_{\max})}(H)}\right\}.
\end{equation}
For $w_j=1/k$, one has $\mathcal B_w|_{\mathcal T}
=\operatorname{Id}/k$, and
$\mathcal E_{M,w,k/N}(H)=\mathcal Q_N(E,F)^2/N$ exactly.
\end{proposition}
\begin{proof}
The variance expression gives
$0\preceq\mathcal B_w\preceq w_{\max}\operatorname{Id}$.
On $\mathcal T$, all coordinates of $w$ positive make
$\mathcal B_w$ positive definite.  Inverse order therefore gives
\[
 \langle H,(\mathcal C_M+N^{-1}\operatorname{Id})^{-1}H\rangle
 \le\mathcal E_{M,w,1/(Nw_{\max})}(H).
\]
Use Theorem~\ref{thm:closedcovariance83}.  The uniform formula
follows by substitution into \eqref{eq:weightedridge84}.
\end{proof}
Theorem~\ref{thm:transportedridge84} compares the same $\tau$ and
\emph{transported} weights.  Replacing $Tw$ by a new uniform vector,
or changing $\tau$ to the new alphabet size divided by $N$, is a
different operation.  Consequently this result does not retroactively
assert arbitrary-processing monotonicity of the original fixed ridge.
Exact evaluation with positive rational weights uses the tangent Gram
system of Proposition~\ref{prop:covarianceexact83}, with the ridge
stiffness $\tau\mathcal B_w$ in place of $G/N$.
